% 5-Qubit Perfect Code Experiment Design
% Optional appendix or future work section

\subsection{Future Work: QEC Code Preparation with Floquet Criteria}
\label{sec:future_floquet}

\subsubsection{Motivation: Testing on Structured Targets}

The experiments in Section~\ref{sec:floquet_results} tested Floquet vs static criteria on \emph{random Haar states}—uniformly distributed targets with no special structure. To assess practical relevance for quantum information tasks, we propose testing on \textbf{quantum error correction (QEC) code states}, which are highly structured and physically motivated targets.

\textbf{Key question:} Does the Floquet advantage persist (or even strengthen) for structured targets that may inherently require higher-body interactions?

\subsubsection{5-Qubit Perfect Code}

The [[5,1,3]] perfect code~\cite{Laflamme_1996} is the smallest code that corrects arbitrary single-qubit errors. It encodes 1 logical qubit into 5 physical qubits with code distance $d_{\text{code}}=3$.

\paragraph{Stabilizer Generators}

The code is defined by four commuting stabilizer operators:
\begin{align}
S_1 &= X Z Z X I \\
S_2 &= I X Z Z X \\
S_3 &= X I X Z Z \\
S_4 &= Z X I X Z
\end{align}

The logical zero state $|0_L\rangle$ is the simultaneous $+1$ eigenstate of all stabilizers:
\begin{equation}
S_j |0_L\rangle = |0_L\rangle \quad \forall j \in \{1,2,3,4\}
\end{equation}

\paragraph{Explicit Construction}

The logical zero state can be written as an equal superposition of 16 basis states with specific signs determined by the stabilizer group structure:
\begin{equation}
|0_L\rangle = \frac{1}{4} \sum_{s \in \mathcal{S}} (-1)^{w(s)} |s\rangle
\end{equation}
where $\mathcal{S}$ is the codespace (16 basis states) and $w(s)$ encodes the sign pattern.

Explicitly, the 16 basis states with their coefficients are:
\begin{align*}
\text{Weight 0:} \quad &+|00000\rangle \\
\text{Weight 2:} \quad &+|10010\rangle, +|01001\rangle, +|10100\rangle, +|01010\rangle \\
\text{Weight 3:} \quad &-|11011\rangle, -|00110\rangle, -|11000\rangle, -|11101\rangle, \\
&-|00011\rangle, -|11110\rangle, -|01111\rangle, -|10001\rangle, \\
&-|01100\rangle, -|10111\rangle \\
\text{Weight 4:} \quad &+|00101\rangle
\end{align*}

\paragraph{Physical Properties}

\begin{itemize}
\item \textbf{Entanglement:} Highly entangled—no product state representation
\item \textbf{Overlap with $|00000\rangle$:} $|\langle 00000 | 0_L \rangle|^2 = 1/16 \approx 0.063$
\item \textbf{Symmetry:} Invariant under stabilizer group (abelian subgroup of Pauli group)
\item \textbf{Distance:} Minimum weight of logical operators is 3 (hence distance-3 code)
\end{itemize}

\subsubsection{Experimental Design}

\paragraph{System Setup}
\begin{itemize}
\item \textbf{Qubits:} $n = 5$ ($d = 2^5 = 32$ dimensional Hilbert space)
\item \textbf{Hamiltonian ensemble:} GEO2LOCAL (1D chain)
  \begin{itemize}
  \item 5 qubits arranged in a line: $1-2-3-4-5$
  \item Nearest-neighbor interactions: $(1,2), (2,3), (3,4), (4,5)$
  \item Total operators: $L = 3n + 9|E| = 15 + 36 = 51$
  \end{itemize}
\item \textbf{Initial state:} $|\psi\rangle = |00000\rangle$
\item \textbf{Target state:} $|\varphi\rangle = |0_L\rangle$ (fixed, deterministic)
\end{itemize}

\paragraph{Protocol}

For each $K \in \{2, 4, 6, 8, 10, 12, 14, 16, 18, 20, 22, 24, 26, 28, 30\}$:
\begin{enumerate}
\item Select $K$ random operators from the GEO2LOCAL ensemble
\item \textbf{Static criterion:} Test whether $\mathbf{Q} + x \mathbf{L} \mathbf{L}^T \succ 0$
  \begin{itemize}
  \item If yes: $|0_L\rangle$ is provably unreachable with these $K$ operators
  \end{itemize}
\item \textbf{Floquet criterion:} Search over 100 random $\boldsymbol{\lambda}$ vectors
  \begin{itemize}
  \item For each $\boldsymbol{\lambda}$: test $\mathbf{Q}_F(\boldsymbol{\lambda}) + x \mathbf{L}_F(\boldsymbol{\lambda}) \mathbf{L}_F(\boldsymbol{\lambda})^T \succ 0$
  \item If any $\boldsymbol{\lambda}$ succeeds: $|0_L\rangle$ is provably unreachable
  \end{itemize}
\item Record success/failure for both criteria
\end{enumerate}

\paragraph{Critical $K$ Values}

Define the \textbf{critical operator count}:
\begin{equation}
K_c^{\text{criterion}} = \max \{ K : \text{criterion proves unreachability} \}
\end{equation}

This is the \emph{largest} $K$ for which the criterion can still detect that $|0_L\rangle$ is unreachable with only $K$ operators.

\textbf{Expected ordering:}
\begin{equation}
K_c^{\text{static}} < K_c^{\text{floquet}} < K_c^{\text{optimal}}
\end{equation}

\paragraph{Predicted Outcomes}

\textbf{Scenario 1: Floquet advantage confirmed}
\begin{itemize}
\item $K_c^{\text{static}} = 6$ (proves unreachable with $K \leq 6$)
\item $K_c^{\text{floquet}} = 10$ (proves unreachable with $K \leq 10$)
\item \textbf{Interpretation:} Floquet criterion is stronger—can detect unreachability with 67\% more operators
\end{itemize}

\textbf{Scenario 2: Both criteria fail early}
\begin{itemize}
\item $K_c^{\text{static}} = 0$, $K_c^{\text{floquet}} = 0$ (neither proves unreachability at any $K$)
\item \textbf{Interpretation:} $|0_L\rangle$ is "generic" enough that moment criteria are insufficient (similar to Haar states)
\end{itemize}

\textbf{Scenario 3: Dramatic Floquet advantage}
\begin{itemize}
\item $K_c^{\text{static}} = 0$, $K_c^{\text{floquet}} = 8$
\item \textbf{Interpretation:} Commutator terms are \emph{essential}—only Floquet can detect unreachability
\end{itemize}

\subsubsection{Physical Interpretation}

\paragraph{Why QEC Codes May Require Commutators}

The stabilizer structure of $|0_L\rangle$ involves \textbf{non-local correlations}—multi-qubit Pauli operators acting simultaneously on distant qubits. These correlations may not be expressible as linear combinations of 2-local operators alone, but could emerge from commutators:
\begin{equation}
[Z_1 Z_2, Z_3 Z_4] \propto Z_1 Z_2 Z_3 Z_4 \quad \text{(4-body term)}
\end{equation}

If $|0_L\rangle$ preparation requires such higher-body terms, the Floquet criterion (which includes commutators) should detect this, while the static criterion (which does not) will fail.

\paragraph{Magnus Order Classification}

A byproduct of this experiment would be a \textbf{classification scheme for QEC codes}:
\begin{itemize}
\item \textbf{Moment-easy codes:} Static criterion succeeds (no commutators needed)
\item \textbf{Floquet-1 codes:} Require $\hat{H}_F^{(1)}$ only (time-averaged Hamiltonian)
\item \textbf{Floquet-2 codes:} Require $\hat{H}_F^{(2)}$ (single commutators, like [[5,1,3]]?)
\item \textbf{Floquet-3 codes:} Require $\hat{H}_F^{(3)}$ (nested commutators)
\item \textbf{Floquet-hard codes:} No Magnus order suffices (e.g., toric code?)
\end{itemize}

This mirrors the GKP code example (lines 654-682 in longer reference), where first-order Magnus suffices for preparation.

\subsubsection{Computational Cost}

\begin{itemize}
\item \textbf{System size:} $d=32$ (vs $d=16$ in main experiment)
\item \textbf{Per-trial cost:} ~1-2 seconds (static), ~100-200 seconds (Floquet with 100 $\lambda$ searches)
\item \textbf{Number of $K$ values:} 15
\item \textbf{Total runtime:} ~50 minutes per criterion
\item \textbf{Grand total:} ~1.5-2 hours for complete experiment
\end{itemize}

Significantly faster than the 11-hour Haar state experiment due to testing a single deterministic target rather than 100 random targets per $K$.

\subsubsection{Implementation Status}

The experiment is \textbf{fully implemented} in \texttt{scripts/run\_5qubit\_code\_experiment.py} with:
\begin{itemize}
\item Explicit construction of $|0_L\rangle$ with correct sign pattern
\item Stabilizer verification (checks all four $S_j$ eigenvalues = +1)
\item Automated search over $K$ values
\item Extraction of critical $K_c$ values
\item Results saved as pickle + JSON summary
\end{itemize}

\textbf{Usage:}
\begin{verbatim}
python3 scripts/run_5qubit_code_experiment.py \
  --n-lambda-search 100 \
  --seed 42
\end{verbatim}

\subsubsection{Extensions}

\paragraph{Multiple Code States}

Compare different logical states of the same code:
\begin{itemize}
\item $|0_L\rangle$ vs $|1_L\rangle$ (related by $X_L = XXXXX$)
\item $|+_L\rangle = (|0_L\rangle + |1_L\rangle)/\sqrt{2}$ (logical superposition)
\end{itemize}

\textbf{Question:} Do certain logical states require higher Magnus order?

\paragraph{Different Codes}

Test the hierarchy across codes:
\begin{itemize}
\item [[5,1,3]] (5-qubit perfect code)
\item [[7,1,3]] (Steane code)
\item [[9,1,3]] (Shor code)
\end{itemize}

\textbf{Question:} Does code size or structure correlate with required Magnus order?

\paragraph{Stabilizer-Based Hamiltonians}

Instead of random GEO2LOCAL operators, use Hamiltonians constructed from stabilizer generators:
\begin{equation}
H_k \in \{ S_1, S_2, S_3, S_4, S_1 S_2, S_1 S_3, \ldots \}
\end{equation}

This is more realistic for QEC implementations where stabilizers are measured and used for error correction.

\subsubsection{Expected Impact}

Success of this experiment would:
\begin{enumerate}
\item Demonstrate Floquet advantage on \textbf{physically relevant targets} (not just random states)
\item Provide a \textbf{classification tool} for assessing QEC code complexity
\item Guide \textbf{experimental design} for code state preparation (which operators + driving schemes are needed?)
\item Connect to existing literature on Floquet codes (e.g., GKP in \cite{Kolesnikow_2024})
\end{enumerate}

This would elevate the Floquet moment criterion from a theoretical curiosity to a \textbf{practical tool for quantum control design}.
